\frfilename{mt226.tex}
\versiondate{16.11.13}
\copyrightdate{2000}

\def\chaptername{Teorema Dasar Kalkulus}
\def\sectionname{Dekomposisi Lebesgue suatu fungsi bervariasi
terbatas}

\newsection{226}

Saya menutup bab ini dengan beberapa catatan mengenai suatu metode untuk
menganalisis fungsi umum yang bervariasi terbatas. Metode ini dapat membantu
memberikan gambaran tentang seperti apa fungsi-fungsi tersebut, walaupun
(selain 226A) hampir tidak diperlukan dalam jilid ini.

\leader{226A}{Jumlah atas himpunan indeks sembarang} \cmmnt{Untuk
memperoleh gambaran lengkap mengenai bagian analisis real ini, sedikit persiapan
akan berguna. Persiapan ini menyangkut konsep jumlah atas himpunan
indeks sembarang, yang sejauh ini cenderung saya hindari.

\header{226Aa}}{\bf (a)} Jika $I$ sembarang himpunan dan
$\langle a_i\rangle_{i\in I}$ sembarang keluarga dalam $[0,\infty]$, kita
tetapkan

\Centerline{$\sum_{i\in I}a_i
=\sup\{\sum_{i\in K}a_i:K\text{ suatu subhimpunan berhingga dari }I\}$,}

\noindent dengan kaidah bahwa $\sum_{i\in\emptyset}a_i=0$.
\cmmnt{(Lihat 112Bd, 222Ba.)} Untuk $a_i\in[-\infty,\infty]$ umum, kita
dapat menetapkan

\Centerline{$\sum_{i\in I}a_i=\sum_{i\in I}a_i^+-\sum_{i\in I}a_i^-$}

\noindent jika ruas ini terdefinisi dalam $[-\infty,\infty]$\cmmnt{, yakni
sekurang-kurangnya salah satu dari $\sum_{i\in I}a_i^+$,
$\sum_{i\in I}a_i^-$ berhingga}, dengan $a^+=\max(a,0)$ dan
$a^-=\max(-a,0)$ untuk setiap $a$. Jika $\sum_{i\in I}a_i$ terdefinisi
dan berhingga, kita mengatakan bahwa $\familyiI{a_i}$ {\bf dapat
dijumlahkan}.

\header{226Ab}{\bf (b)}\cmmnt{ Karena buku ini membahas teori ukuran,
saya akan langsung menguraikan hubungan antara dapat dijumlahkannya keluarga
semacam ini dan konsep integrasi yang sesuai.} Untuk setiap himpunan $I$,
kita mempunyai `ukuran cacah' $\mu$ pada $I$ yang bersesuaian\cmmnt{
(112Bd)}. Setiap\cmmnt{ subhimpunan dari $I$ terukur, sehingga setiap}
keluarga $\langle a_i\rangle_{i\in I}$ bilangan real merupakan fungsi
bernilai real terukur pada $I$. Suatu\cmmnt{ subhimpunan dari $I$
berukuran berhingga jika dan hanya jika subhimpunan itu berhingga; jadi}
fungsi bernilai real $f$ pada $I$ disebut `sederhana' jika
$K=\{i:f(i)\ne 0\}$ berhingga. \cmmnt{Dalam hal ini,

\Centerline{$\int fd\mu=\sum_{i\in K}f(i)=\sum_{i\in I}f(i)$}

\noindent sebagaimana didefinisikan dalam bagian (a). Ukuran $\mu$
semifinit (211Nc), sehingga fungsi nonnegatif $f$ terintegralkan jika dan
hanya jika $\int f=\sup_{\mu K<\infty}\int_Kf$ berhingga (213B); tetapi
tentu saja supremum ini tidak lain adalah

\Centerline{$\sup\{\sum_{i\in K}f(i):K\subseteq I\text{ berhingga}\}
=\sum_{i\in I}f(i)$.}


\noindent}%end of comment
Sekarang suatu fungsi sembarang $f:I\to\Bbb R$ terintegralkan\cmmnt{ jika dan
hanya jika fungsi itu terukur dan $\int|f|d\mu<\infty$, yakni,} jika dan
hanya jika $\sum_{i\in I}|f(i)|<\infty$, dan dalam hal ini

\dvro{\Centerline{$\int fd\mu=\sum_{i\in I}f(i)$,}}
{\Centerline{$\int fd\mu
=\int f^+d\mu-\int f^-d\mu
=\sum_{i\in I}f(i)^+-\sum_{i\in I}f(i)^-
=\sum_{i\in I}f(i)$,}}

\noindent\cmmnt{dengan menuliskan $f^{\pm}(i)=f(i)^{\pm}$ untuk setiap
$i$. }Jadi kita mempunyai

\Centerline{$\sum_{i\in I}a_i=\int_Ia_i\mu(di)$,}

\noindent dan aturan baku yang memperbolehkan $\infty$ sebagai nilai suatu
integral\cmmnt{ (133A, 135F)} sangat selaras\cmmnt{ dengan} penafsiran
dalam (a) di atas.

\header{226Ac}{\bf (c)}\cmmnt{ Dengan demikian, dan seperti yang dapat
diduga, operasi penjumlahan merupakan operasi linear pada ruang linear
keluarga bilangan real yang dapat dijumlahkan.

} Saya mencatat di sini bahwa konsep keterjumlahan ini bersifat
`mutlak'; suatu keluarga $\langle a_i\rangle_{i\in I}$ dapat dijumlahkan
jika dan hanya jika keluarga itu dapat dijumlahkan mutlak.
\cmmnt{Hal ini diperlukan karena penjumlahannya juga harus
`tak bersyarat'; kita tidak mempunyai struktur pada himpunan sembarang $I$
yang dapat mengarahkan kita untuk mengambil jumlah dalam suatu urutan tertentu.
Lihat 226Xa. Secara khusus, saya membedakan
`$\sum_{n\in\Bbb N}a_n$', yang dalam buku ini selalu akan ditafsirkan
seperti dalam 226A di atas, dari `$\sum_{n=0}^{\infty}a_n$', yang
(jika pembedaan ini berpengaruh) harus ditafsirkan sebagai
$\lim_{m\to\infty}\sum_{n=0}^ma_n$. Jadi, misalnya,
$\sum_{n=0}^{\infty}\Bover{(-1)^n}{n+1}=\ln 2$, sedangkan
$\sum_{n\in\Bbb N}\Bover{(-1)^n}{n+1}$ tidak terdefinisi. Tentu saja
$\sum_{n=0}^{\infty}a_n=\sum_{n\in\Bbb N}a_n$ kapan pun jumlah yang
terakhir terdefinisi dalam $[-\infty,\infty]$.}

\header{226Ad}{\bf (d)}\cmmnt{ Ada pendekatan lain yang sangat penting
terhadap jumlah yang diuraikan di sini.} Jika
$\langle a_i\rangle_{i\in I}$ merupakan suatu keluarga bilangan real yang
(secara mutlak) dapat dijumlahkan, maka untuk setiap $\epsilon>0$ terdapat suatu
$K\subseteq I$ berhingga sedemikian sehingga
$\sum_{i\in I\setminus K}|a_i|\le\epsilon$. \prooflet{\Prf\ Ini tidak
lain daripada kasus khusus 225A; terdapat suatu himpunan $K$ dengan
$\mu K<\infty$ sedemikian sehingga
$\int_{I\setminus K}|a_i|\mu(di)\le\epsilon$, tetapi

\Centerline{$\int_{I\setminus K}|a_i|\mu(di)
=\sum_{i\in I\setminus K}|a_i|$. \Qed}

\noindent}\cmmnt{(Tentu saja ada pembuktian `langsung' atas hasil ini
dari definisi dalam (a), tanpa menyebut ukuran atau integral. Namun saya
rasa Anda akan melihat bahwa pembuktian tersebut bertumpu pada gagasan yang
sama dengan pembuktian 225A.) }Akibatnya, untuk sembarang keluarga
$\langle a_i\rangle_{i\in I}$ bilangan real dan sembarang $s\in\Bbb R$,
pernyataan-pernyataan berikut ekuivalen:

\inset{(i) $\sum_{i\in I}a_i=s$;}

\inset{(ii) untuk setiap $\epsilon>0$ terdapat suatu $K\subseteq I$ berhingga
sedemikian sehingga $|s-\sum_{i\in J}a_i|\le\epsilon$ kapan pun $J$
berhingga dan $K\subseteq J\subseteq I$.}

\noindent\prooflet{\Prf\ {\bf (i)$\Rightarrow$(ii)} Ambil $K$ sedemikian
sehingga $\sum_{i\in
I\setminus K}|a_i|\le\epsilon$. Jika $K\subseteq J\subseteq I$, maka

\Centerline{$|s-\sum_{i\in J}a_i|=|\sum_{i\in I\setminus J}a_i|
\le\sum_{i\in I\setminus K}|a_i|\le\epsilon$.}

{\bf (ii)$\Rightarrow$(i)} Misalkan $\epsilon>0$, dan misalkan
$K\subseteq I$ seperti yang diuraikan dalam (ii). Jika
$J\subseteq I\setminus K$ sembarang himpunan berhingga, tetapkan
$J_1=\{i:i\in J,\,a_i\ge 0\}$, $J_2=J\setminus J_1$. Kita mempunyai

$$\eqalign{\sum_{i\in J}|a_i|
&=|\sum_{i\in J_1\cup K}a_i-\sum_{i\in J_2\cup K}a_i|\cr
&\le|s-\sum_{i\in J_1\cup K}a_i|+|s-\sum_{i\in J_2\cup K}a_i|
\le 2\epsilon.\cr}$$

\noindent Karena $J$ sembarang, $\sum_{i\in I\setminus K}|a_i|\le
2\epsilon$ dan

\Centerline{$\sum_{i\in I}|a_i|\le\sum_{i\in K}|a_i|+2\epsilon<\infty$.}

\noindent Oleh karena itu, $\sum_{i\in I}a_i$ terdefinisi dengan baik dalam
$\Bbb R$. Selain itu,

\Centerline{$|s-\sum_{i\in I}a_i|
\le|s-\sum_{i\in K}a_i|+|\sum_{i\in I\setminus K}a_i|
\le\epsilon+\sum_{i\in I\setminus K}|a_i|
\le 3\epsilon$.}

\noindent Karena $\epsilon$ sembarang, $\sum_{i\in I}a_i=s$, sebagaimana
diperlukan. \Qed}

\cmmnt{Dengan cara ini kita menyatakan $\sum_{i\in I}a_i$ secara langsung
sebagai suatu limit; kita dapat menuliskannya sebagai

\Centerline{$\sum_{i\in I}a_i=\lim_{K\uparrow I}\sum_{i\in K}a_i$,}

\noindent dengan pengertian bahwa dalam rumus di ruas kanan kita hanya
memperhatikan himpunan berhingga $K$.
}%end of comment

\header{226Ae}{\bf (e)}\cmmnt{ Pendekatan lain lagi adalah melalui fakta
berikut.} Jika $\sum_{i\in I}|a_i|<\infty$, maka\cmmnt{ untuk sembarang
$\delta>0$ himpunan $\{i:|a_i|\ge\delta\}$ berhingga, bahkan dapat
mempunyai paling banyak ${1\over{\delta}}\sum_{i\in I}|a_i|$ anggota;
akibatnya}

\Centerline{$J=\{i:a_i\ne 0\}
=\bigcup_{n\in\Bbb N}\{i:|a_i|\ge 2^{-n}\}$}

\noindent terhitung\cmmnt{ (1A1F)}. Jika $J$ berhingga, maka\cmmnt{ tentu
saja} $\sum_{i\in I}a_i=\sum_{i\in J}a_i$ menjadi suatu jumlah berhingga.
Jika tidak, kita dapat mengenumerasi $J$ sebagai $\sequencen{j_n}$, dan
kita akan mempunyai

\Centerline{$\sum_{i\in I}a_i=\sum_{i\in
J}a_i=\lim_{n\to\infty}\sum_{k=0}^na_{j_k}=\sum_{n=0}^{\infty}a_{j_n}$
\dvro{.}{}}

\noindent\cmmnt{(menggunakan (d) untuk mereduksi jumlah
$\sum_{i\in J}a_i$ menjadi limit jumlah-jumlah berhingga). }Sebaliknya,
jika $\langle a_i\rangle_{i\in I}$ sedemikian sehingga terdapat suatu
$J=\{i:a_i\ne 0\}$ tak hingga terhitung yang dienumerasi sebagai
$\sequencen{j_n}$, dan jika $\sum_{n=0}^{\infty}|a_{j_n}|<\infty$, maka
$\sum_{i\in I}a_i$ akan sama dengan $\sum_{n=0}^{\infty}a_{j_n}$.

\spheader 226Af
%new 2008
\cmmnt{Kelak akan berguna untuk mempunyai suatu fragmen teori umum.}
Misalkan $I$ dan $J$ adalah himpunan, dan
$\langle a_{ij}\rangle_{i\in I,j\in J}$ suatu keluarga dalam
$[0,\infty]$. Maka

\Centerline{$\sum_{(i,j)\in I\times J}a_{ij}
=\sum_{i\in I}(\sum_{j\in J}a_{ij})
=\sum_{j\in J}(\sum_{i\in I}a_{ij})$.}

\prooflet{\noindent\Prf\ (i) Jika
$\sum_{(i,j)\in I\times J}a_{ij}>u$, maka terdapat suatu himpunan
$M\subseteq I\times J$ berhingga sedemikian sehingga
$\sum_{(i,j)\in M}a_{ij}>u$. Sekarang $K=\{i:\exists j\ (i,j)\in M\}$ dan
$L=\{j:\exists i\ (i,j)\in M\}$ berhingga, sehingga

$$\eqalignno{\sum_{i\in I}\sum_{j\in J}a_{ij}
&\ge\sum_{i\in K}\sum_{j\in J}a_{ij}
\ge\sum_{i\in K}\sum_{j\in L}a_{ij}\cr
\displaycause{karena $\sum_{j\in J}a_{ij}\ge\sum_{j\in L}a_{ij}$ untuk setiap
$i$}
&=\sum_{(i,j)\in K\times L}a_{ij}
\ge\sum_{(i,j)\in M}a_{ij}
>u.\cr}$$

\noindent Karena $u$ sembarang,
$\sum_{i\in I}\sum_{j\in J}a_{ij}\ge\sum_{(i,j)\in I\times J}a_{ij}$.
(ii) Jika $\sum_{i\in I}\sum_{j\in J}a_{ij}>u\ge 0$, terdapat suatu
himpunan $K\subseteq I$ tak kosong dan berhingga sedemikian sehingga
$\sum_{i\in K}\sum_{j\in J}a_{ij}>u$.
Misalkan $\epsilon\in\ooint{0,1}$ sedemikian sehingga
$\sum_{i\in K}\sum_{j\in J}a_{ij}>u+\epsilon$, dan tetapkan
$\delta=\Bover{\epsilon}{\#(K)}$. Untuk setiap $i\in K$, tetapkan
$\gamma_i=\min(u+1,\sum_{j\in J}a_{ij})-\delta$;
maka

\Centerline{$\epsilon+\sum_{i\in K}\gamma_i
=\sum_{i\in K}\min(u+1,\sum_{j\in J}a_{ij})
\ge\min(u+1,\sum_{i\in K}\sum_{j\in J}a_{ij})
>u+\epsilon$,}

\noindent sehingga $\sum_{i\in K}\gamma_i>u$. Untuk setiap $i\in K$,
$\gamma_i<\sum_{j\in J}a_{ij}$, sehingga terdapat suatu
$L_i\subseteq J$ berhingga sedemikian sehingga
$\sum_{j\in L_i}a_{ij}\ge\gamma_i$.
Tetapkan $M=\{(i,j):i\in K$, $j\in L_i\}$, sehingga $M$ merupakan
subhimpunan berhingga dari $I\times J$; maka

\Centerline{$\sum_{(i,j)\in I\times J}a_{ij}
\ge\sum_{(i,j)\in M}a_{ij}
=\sum_{i\in K}\sum_{j\in L_i}a_{ij}
\ge\sum_{i\in K}\gamma_i
>u$.}

\noindent Karena $u$ sembarang,
$\sum_{(i,j)\in I\times J}a_{ij}\ge\sum_{i\in I}\sum_{j\in J}a_{ij}$ dan
kedua jumlah ini sama. (iii) Demikian pula,
$\sum_{(i,j)\in I\times J}a_{ij}=\sum_{j\in J}\sum_{i\in I}a_{ij}$.\ \Qed}

\leader{226B}{Fungsi saltus}\cmmnt{ Sekarang kita siap membahas suatu
jenis khusus fungsi bervariasi terbatas pada $\Bbb R$.} Misalkan $a<b$
dalam $\Bbb R$.

\header{226Ba}{\bf (a)} Suatu {\bf fungsi saltus} (bernilai real) pada $[a,b]$
adalah fungsi $F:[a,b]\to\Bbb R$ yang dapat dinyatakan dalam bentuk

\Centerline{$F(x)=\sum_{t\in\coint{a,x}}u_t+\sum_{t\in[a,x]}v_t$}

\noindent untuk $x\in [a,b]$, dengan $\langle
u_t\rangle_{t\in\coint{a,b}}$ dan
$\langle v_t\rangle_{t\in[a,b]}$ keluarga-keluarga bernilai real
sedemikian sehingga $\sum_{t\in\coint{a,b}}|u_t|$ dan
$\sum_{t\in[a,b]}|v_t|$ berhingga.

\header{226Bb}{\bf (b)} Untuk setiap fungsi $F:[a,b]\to\Bbb R$, kita
dapat menuliskan

\Centerline{$F(x^+)=\lim_{y\downarrow x}F(y)$ jika $x\in\coint{a,b}$ dan
limit tersebut ada,}

\Centerline{$F(x^-)=\lim_{y\uparrow x}F(y)$ jika $x\in\ocint{a,b}$ dan
limit tersebut ada.}

\noindent\cmmnt{(Saya harap hal ini tidak menimbulkan kebingungan dengan
penafsiran lain atas $x^+$ sebagai $\max(x,0)$.) }Perhatikan bahwa jika
$F$ suatu fungsi saltus, sebagaimana didefinisikan dalam (a), dengan
keluarga-keluarga terkait $\langle u_t\rangle_{t\in\coint{a,b}}$ dan
$\langle v_t\rangle_{t\in[a,b]}$, maka $v_a=F(a)$,
$v_x=F(x)-F(x^-)$ untuk $x\in\ocint{a,b}$ dan
$u_x=F(x^+)-F(x)$ untuk $x\in\coint{a,b}$.
\prooflet{\Prf\ Misalkan
$\epsilon>0$. Sebagaimana dicatat dalam 226Ad, terdapat suatu
$K\subseteq[a,b]$ berhingga sedemikian sehingga

\Centerline{$\sum_{t\in\coint{a,b}\setminus
K}|u_t|+\sum_{t\in[a,b]\setminus K}|v_t|\le\epsilon$.}

\noindent Untuk $x\in[a,b]$ yang diberikan, misalkan $\delta>0$ sedemikian
sehingga $[x-\delta,x+\delta]$ tidak memuat titik $K$ selain mungkin $x$.
Dalam hal ini, jika $\max(a,x-\delta)\le y<x$, kita pasti mempunyai

$$\eqalign{|F(y)-(F(x)-v_x)|
&=|\sum_{t\in\coint{y,x}}u_t+\sum_{t\in\ooint{y,x}}v_t|\cr
&\le\sum_{t\in\coint{a,b}\setminus K}|u_t|
  +\sum_{t\in[a,b]\setminus K}|v_t|
\le\epsilon,\cr}$$

\noindent sedangkan jika $x<y\le\min(b,x+\delta)$, kita akan mempunyai

$$\eqalign{|F(y)-(F(x)+u_x)|
&=|\sum_{t\in\ooint{x,y}}u_t+\sum_{t\in\ocint{x,y}}v_t|\cr
&\le\sum_{t\in\coint{a,b}\setminus K}|u_t|
  +\sum_{t\in[a,b]\setminus K}|v_t|
\le\epsilon.\cr}$$

\noindent Karena $\epsilon$ sembarang, kita memperoleh
$F(x^-)=F(x)-v_x$ (jika $x>a$) dan $F(x^+)=F(x)+u_x$ (jika $x<b$). \Qed}

\cmmnt{Akibatnya, }$F$ kontinu di $x\in\ooint{a,b}$
jika dan hanya jika $u_x=v_x=0$,
sedangkan $F$ kontinu di $a$ jika dan hanya jika $u_a=0$, dan $F$ kontinu
di $b$ jika dan hanya jika $v_b=0$. Secara khusus,
$\{x:x\in[a,b],\,F\text{ tidak kontinu di }x\}$ terhitung\cmmnt{
(lihat 226Ae)}.

\header{226Bc}{\bf (c)} Jika $F$ suatu fungsi saltus yang didefinisikan
pada $[a,b]$, dengan keluarga-keluarga terkait
$\langle u_t\rangle_{t\in\coint{a,b}}$ dan
$\langle v_t\rangle_{t\in[a,b]}$, maka $F$ bervariasi terbatas pada
$[a,b]$, dan

\Centerline{$\Var_{[a,b]}(F)
\le\sum_{t\in\coint{a,b}}|u_t|+\sum_{t\in\ocint{a,b}}|v_t|$.}

\noindent\prooflet{\Prf\ Jika $a\le x<y\le b$, maka

\Centerline{$F(y)-F(x)=u_x+\sum_{t\in\ooint{x,y}}(u_t+v_t)+v_y$,}

\noindent sehingga

\Centerline{$|F(y)-F(x)|
\le\sum_{t\in\coint{x,y}}|u_t|+\sum_{t\in\ocint{x,y}}|v_t|$.}

\noindent Jika $a\le a_0\le a_1\le\ldots\le a_n\le b$, maka

$$\eqalign{\sum_{i=1}^n|F(a_i)-F(a_{i-1})|
&\le\sum_{i=1}^n\bigl(\sum_{t\in\coint{a_{i-1},a_i}}|u_t|
   +\sum_{t\in\ocint{a_{i-1},a_i}}|v_t|\bigr)\cr
&\le\sum_{t\in\coint{a,b}}|u_t|+\sum_{t\in\ocint{a,b}}|v_t|.\cr}$$

\noindent Akibatnya,

\Centerline{$\Var_{[a,b]}(F)
\le\sum_{t\in\coint{a,b}}|u_t|+\sum_{t\in\ocint{a,b}}|v_t|<\infty$. \Qed}}

\header{226Bd}{\bf (d)} Ketaksamaan dalam (c) sebenarnya merupakan
kesamaan. \prooflet{Untuk melihatnya, perhatikan terlebih dahulu bahwa
jika $a\le x<y\le b$, maka
$\Var_{[x,y]}(F)\ge|u_x|+|v_y|$. \Prf\ Saya mencatat dalam (b) bahwa
$u_x=\lim_{t\downarrow x}F(t)-F(x)$ dan
$v_y=F(y)-\lim_{t\uparrow y}F(t)$.
Jadi, untuk $\epsilon>0$ yang diberikan, kita dapat menemukan $t_1$,
$t_2$ sedemikian sehingga $x<t_1\le t_2<y$ dan

\Centerline{$|F(t_1)-F(x)|\ge|u_x|-\epsilon$,
\quad$|F(y)-F(t_2)|\ge|v_y|-\epsilon$.}

\noindent Sekarang

\Centerline{$\Var_{[x,y]}(F)
\ge|F(t_1)-F(x)|+|F(t_2)-F(t_1)|+|F(y)-F(t_2)|
\ge|u_x|+|v_y|-2\epsilon$.}

\noindent Karena $\epsilon$ sembarang, kita memperoleh hasil tersebut.
\Qed

Sekarang, untuk $a\le t_0<t_1<\ldots<t_n\le b$ yang diberikan, kita pasti
mempunyai

$$\eqalignno{\Var_{[a,b]}(F)
&\ge\sum_{i=1}^n\Var_{[t_{i-1},t_i]}(F)\cr
\noalign{\noindent(menggunakan 224Cc)}
&\ge\sum_{i=1}^n|u_{t_{i-1}}|+|v_{t_i}|.\cr}$$

\noindent Karena $t_0,\ldots,t_n$ sembarang,

\Centerline{$\Var_{[a,b]}(F)
\ge\sum_{t\in\coint{a,b}}|u_t|+\sum_{t\in\ocint{a,b}}|v_t|$,}

\noindent sebagaimana diperlukan.}

\header{226Be}{\bf (e)} Karena suatu fungsi saltus bervariasi
terbatas\cmmnt{ ((c) di atas)}, fungsi itu terdiferensialkan hampir di
mana-mana\cmmnt{ (224I)}.
Sebenarnya turunannya nol hampir di mana-mana. \prooflet{\Prf\ Misalkan
$F:[a,b]\to\Bbb R$ suatu fungsi saltus, dengan keluarga-keluarga terkait
$\langle u_t\rangle_{t\in\coint{a,b}}$ dan
$\langle v_t\rangle_{t\in[a,b]}$.
Misalkan $\epsilon>0$. Misalkan $K\subseteq[a,b]$ suatu himpunan berhingga
sedemikian sehingga

\Centerline{$\sum_{t\in\coint{a,b}\setminus
K}|u_t|+\sum_{t\in[a,b]\setminus K}|v_t|\le\epsilon$.}

\noindent Tetapkan

$$\eqalign{u'_t
&=u_t\text{ jika }t\in\coint{a,b}\cap K,\cr
&=0\text{ jika }t\in\coint{a,b}\setminus K,\cr
v'_t&=v_t\text{ jika }t\in K,\cr
&=0\text{ jika }t\in[a,b]\setminus K,\cr
u''_t&=u_t-u'_t\text{ untuk }t\in\coint{a,b},\cr
v''_t&=v_t-v'_t\text{ untuk }t\in[a,b].\cr}$$

\noindent Misalkan $G$ dan $H$ adalah fungsi-fungsi saltus yang bersesuaian dengan
$\langle u'_t\rangle_{t\in\coint{a,b}}$,
$\langle v'_t\rangle_{t\in[a,b]}$,
$\langle u''_t\rangle_{t\in\coint{a,b}}$ dan
$\langle v''_t\rangle_{t\in[a,b]}$, sehingga $F=G+H$. Maka $G'(t)=0$
untuk setiap $t\in\ooint{a,b}\setminus K$, karena
$\ooint{a,b}\setminus K$ merupakan gabungan berhingga interval-interval
terbuka, dan $G$ konstan pada setiap interval tersebut. Jadi $G'=0$ hampir di
mana-mana dan
$F'\eae H'$. Di sisi lain,

\Centerline{$\int_a^b|H'|\le\Var_{[a,b]}(H)
=\sum_{t\in\coint{a,b}\setminus K}|u_t|
  +\sum_{t\in\ocint{a,b}\setminus K}|v_t|
\le\epsilon$,}

\noindent dengan menggunakan 224I dan (d) di atas. Jadi

\Centerline{$\int_a^b|F'|=\int_a^b|H'|\le\epsilon$.}

\noindent Karena $\epsilon$ sembarang, $\int_a^b|F'|=0$ dan $F'=0$
hampir di mana-mana, sebagaimana diklaim. \Qed}

\leader{226C}{Dekomposisi Lebesgue bagi fungsi bervariasi
terbatas}  Ambil $a$, $b\in\Bbb R$ dengan $a<b$.

\header{226Ca}{\bf (a)} Jika $F:[a,b]\to\Bbb R$ tak-menurun, tetapkan
$v_a=0$, $v_t=F(t)-F(t^-)$ untuk $t\in\ocint{a,b}$, $u_t=F(t^+)-F(t)$ untuk
$t\in\coint{a,b}$\cmmnt{, dengan mendefinisikan $F(t^+)$, $F(t^-)$ seperti dalam 226Bb}.
Maka semua $v_t$, $u_t$
nonnegatif, dan\cmmnt{ jika $a<t_0<t_1<\ldots<t_n< b$, maka

\Centerline{$\sum_{i=0}^n(u_{t_i}+v_{t_i})
=\sum_{i=0}^n(F(t_i^+)-F(t_i^-))\le F(b)-F(a)$.}

\noindent Oleh karena itu,}
$\sum_{t\in\coint{a,b}}u_t$ dan $\sum_{t\in[a,b]}v_t$ keduanya berhingga.
Misalkan $F_p$ adalah fungsi saltus yang bersesuaian\cmmnt{, sebagaimana didefinisikan dalam
226Ba, sehingga

\Centerline{$F_p(x)
=F(a^+)-F(a)+\sum_{t\in\ooint{a,x}}(F(t^+)-F(t^-))+F(x)
-F(x^-)$}

\noindent jika $a<x\le b$}.   \cmmnt{Jika $a\le x<y\le b$, maka

$$\eqalign{F_p(y)-F_p(x)
&=F(x^+)-F(x)+\sum_{t\in\ooint{x,y}}(F(t^+)-F(t^-))+F(y)-F(y^-)
\cr
&\le F(y)-F(x)\cr}$$

\noindent sebab jika $x=t_0<t_1<\ldots<t_n<t_{n+1}=y$, maka

$$\eqalign{F(x^+)-F(x)
&+\sum_{i=1}^{n}(F(t_i^+)-F(t_i^-))+F(y)-F(y^-)\cr
&=F(y)-F(x)-\sum_{i=1}^{n+1}(F(t_i^-)-F(t_{i-1}^+))
\le F(y)-F(x).\cr}$$

\noindent Oleh karena itu, baik }$F_p$ maupun $F_c=F-F_p$
tak-menurun.   \cmmnt{Selain itu, karena

\Centerline{$F_p(a)=0=v_a$,}

\Centerline{$F_p(t)-F_p(t^-)=v_t=F(t)-F(t^-)$ untuk
$t\in\ocint{a,b}$,}

\Centerline{$F_p(t^+)-F_p(t)=u_t=F(t^+)-F(t)$ untuk
$t\in\coint{a,b}$,}

\noindent kita akan mempunyai

\Centerline{$F_c(a)=F(a)$,}

\Centerline{$F_c(t)=F_c(t^-)$ untuk
$t\in\ocint{a,b}$,}

\Centerline{$F_c(t)=F_c(t^+)$ untuk $t\in\coint{a,b}$,}

\noindent dan }$F_c$ kontinu.

Jelaslah bahwa penyajian $F=F_p+F_c$ ini sebagai jumlah suatu fungsi saltus
dan suatu fungsi kontinu adalah unik, kecuali bahwa kita bebas menambahkan
suatu konstanta pada salah satunya jika konstanta itu dikurangkan dari yang lain.

\header{226Cb}{\bf (b)}\cmmnt{ Dengan tetap menganggap $F:[a,b]\to\Bbb R$
tak-menurun, kita tahu bahwa $F'$ terintegralkan (222C); terlebih lagi,
$F'\eae F'_c$, menurut 226Be.}   Tetapkan $F_{ac}(x)=F(a)+\int_a^xF'$
untuk setiap $x\in[a,b]$.   \cmmnt{Kita mempunyai

\Centerline{$F_{ac}(y)-F_{ac}(x)=\int_x^yF_c'\le F_c(y)-F_c(x)$}

\noindent untuk $a\le x\le y\le b$ (sekali lagi menurut 222C), sehingga} $F_{cs}=F_c-F_{ac}$
masih tak-menurun;  \cmmnt{$F_{ac}$ kontinu (225A),
sehingga }$F_{cs}$
kontinu;  \cmmnt{$F_{ac}'\eae F'\eae F_c'$ (222E),
sehingga }$F_{cs}'=0$ hampir di mana-mana.

Sekali lagi, penyajian $F_c=F_{ac}+F_{cs}$ sebagai jumlah suatu fungsi kontinu
mutlak dan suatu fungsi dengan turunan nol hampir di mana-mana adalah unik,
kecuali untuk kemungkinan memindahkan suatu konstanta dari yang satu ke yang
lain\cmmnt{, karena dua fungsi kontinu mutlak yang
turunannya sama hampir di mana-mana hanya dapat berbeda dengan suatu konstanta
(225D)}.

\header{226Cc}{\bf (c)} Dengan merangkum\cmmnt{ semua} hal ini:  jika
$F:[a,b]\to\Bbb R$ adalah sembarang
fungsi tak-menurun, fungsi tersebut dapat diuraikan sebagai $F_p+F_{ac}+F_{cs}$, dengan
$F_p$ fungsi saltus, $F_{ac}$ kontinu mutlak, dan
$F_{cs}$ kontinu dan turunannya ada serta bernilai nol hampir di
mana-mana;  ketiga komponennya tak-menurun;  dan penguraian itu
unik jika kita menetapkan bahwa $F_{ac}(a)=F(a)$ dan
$F_p(a)=F_{cs}(a)=0$.

Fungsi Cantor $f:[0,1]\to[0,1]$\cmmnt{ (134H)} kontinu dan
$f'=0$ hampir di mana-mana\cmmnt{ (134Hb)}, sehingga $f_p=f_{ac}=0$ dan $f=f_{cs}$.
Dengan menetapkan $g(x)=\bover12(x+f(x))$ untuk $x\in[0,1]$\cmmnt{, seperti dalam 134I},
kita memperoleh $g_p(x)=0$, $g_{ac}(x)={x\over 2}$ dan
$g_{cs}(x)={1\over 2}f(x)$.

\header{226Cd}{\bf (d)} Sekarang andaikan bahwa $F:[a,b]\to\Bbb R$
bervariasi terbatas.
Maka fungsi tersebut dapat dinyatakan sebagai selisih $G-H$ antara
fungsi-fungsi tak-menurun\cmmnt{ (224D)}.   Jadi, dengan menuliskan $F_p=G_p-H_p$,
dan seterusnya, kita dapat menyatakan $F$ sebagai jumlah $F_p+F_{cs}+F_{ac}$, dengan $F_p$
fungsi saltus, $F_{ac}$ kontinu mutlak, $F_{cs}$
kontinu, $F_{cs}'=0$ hampir di mana-mana, $F_{ac}(a)=F(a)$ dan $F_{cs}(a)=F_{p}(a)=0$.
Di bawah syarat-syarat ini penyajian tersebut unik\dvro{.}{, karena (sebagai
contoh) $F_p(t^+)-F_p(t)=F(t^+)-F(t)$ untuk $t\in\coint{a,b}$, sedangkan
$F_{ac}'\eae(F-F_p)'\eae F'$.}

Ini adalah suatu {\bf dekomposisi Lebesgue} dari fungsi $F$.
\cmmnt{(Saya harus mengatakan `suatu' dekomposisi Lebesgue karena tentu saja
penetapan $F_{ac}(a)=F(a)$, $F_{p}(a)=F_{cs}(a)=0$ bersifat sebarang.)}   Saya akan menyebut $F_p$ sebagai {\bf bagian saltus} dari $F$.

\leader{226D}{Fungsi kompleks} \cmmnt{Modifikasi yang diperlukan untuk
menangani fungsi kompleks bersifat elementer.

\medskip

}{\bf (a)} Jika $I$ sembarang himpunan dan
$\langle a_j\rangle_{j\in I}$ suatu keluarga bilangan kompleks, maka
pernyataan-pernyataan berikut ekuivalen:

\inset{(i) $\sum_{j\in I}|a_j|<\infty$;}

\inset{(ii) terdapat $s\in\Bbb C$ sedemikian sehingga untuk setiap $\epsilon>0$
terdapat suatu himpunan berhingga $K\subseteq I$ sedemikian sehingga
$|s-\sum_{j\in J}a_j|\le\epsilon$ kapan pun $J$ berhingga dan
$K\subseteq J\subseteq I$.}

\noindent Dalam hal ini

\Centerline{$s=\sum_{j\in I}\Real(a_j)+i\sum_{j\in
I}\Imag(a_j)=\int_Ia_j\mu(dj)$,}

\noindent dengan $\mu$ ukuran cacah pada $I$, dan kita menulis
$s=\sum_{j\in I}a_j$.

\header{226Db}{\bf (b)} Jika $a<b$ dalam $\Bbb R$, suatu {\bf fungsi
saltus} kompleks pada $[a,b]$ adalah fungsi $F:[a,b]\to\Bbb C$ yang dapat dinyatakan dalam
bentuk

\Centerline{$F(x)=\sum_{t\in\coint{a,x}}u_t+\sum_{t\in[a,x]}v_t$}

\noindent untuk $x\in [a,b]$, dengan $\langle
u_t\rangle_{t\in\coint{a,b}}$,
$\langle v_t\rangle_{t\in[a,b]}$ keluarga-keluarga bernilai kompleks sedemikian sehingga
$\sum_{t\in\coint{a,b}}|u_t|$ dan
$\sum_{t\in[a,b]}|v_t|$ berhingga\cmmnt{;  yakni, jika bagian real dan
bagian imajiner dari $F$ adalah fungsi saltus}.   Dalam hal ini $F$
kontinu kecuali pada terhitung banyak titik, dan turunannya ada serta bernilai
nol hampir di mana-mana dalam $[a,b]$, serta

\Centerline{$u_x=\lim_{t\downarrow x}F(t)-F(x)$ untuk setiap
$x\in\coint{a,b}$,}

\Centerline{$v_x=\lim_{t\uparrow x}F(x)-F(t)$ untuk setiap
$x\in\ocint{a,b}$\dvro{.}{}}

\noindent\cmmnt{(terapkan hasil-hasil 226B pada bagian real dan imajiner
dari $F$).   }$F$ bervariasi terbatas, dan variasinya adalah

\Centerline{$\Var_{[a,b]}(F)
=\sum_{t\in\coint{a,b}}|u_t|+\sum_{t\in\ocint{a,b}}|v_t|$\dvro{.}{}}

\cmmnt{\noindent (ulangi argumen-argumen 226Bc-d).}

\header{226Dc}{\bf (c)} Jika $F:[a,b]\to\Bbb C$ adalah fungsi bervariasi
terbatas, dengan $a<b$ dalam $\Bbb R$, maka fungsi tersebut secara unik dapat dinyatakan sebagai
$F=F_p+F_{cs}+F_{ac}$, dengan $F_p$ fungsi saltus, $F_{ac}$
kontinu mutlak, $F_{cs}$ kontinu dan turunannya bernilai nol
hampir di mana-mana, serta $F_{ac}(a)=F(a)$, $F_p(a)=F_{cs}(a)=0$.
\cmmnt{(Terapkan
226C pada bagian real dan imajiner dari $F$.)}

\leader{226E}{}\cmmnt{ Sebagai suatu latihan elementer dalam bahasa
226A, saya sisipkan sebuah versi dari suatu teorema B. Levi
%Levi 1906
yang kadang-kadang berguna.

\medskip

\noindent}{\bf Proposisi} Misalkan $(X,\Sigma,\mu)$ suatu ruang ukur, $I$
suatu himpunan {\it terhitung}, dan $\familyiI{f_i}$ suatu keluarga fungsi
bernilai real atau kompleks yang $\mu$-terintegralkan sedemikian sehingga $\sum_{i\in I}\int|f_i|d\mu$
berhingga.   Maka
$f(x)=\sum_{i\in I}f_i(x)$ terdefinisi hampir di mana-mana dan
$\int fd\mu=\sum_{i\in I}\int f_id\mu$.

\proof{ Jika $I$ berhingga, hal ini elementer.   Jika tidak, karena harus
terdapat suatu bijeksi antara $I$ dan $\Bbb N$, kita boleh menganggap
$I=\Bbb N$.   Dengan menetapkan $g_n=\sum_{i=0}^n|f_i|$ untuk setiap $n$, kita mempunyai
suatu barisan tak-menurun $\sequencen{g_n}$ dari fungsi-fungsi terintegralkan sedemikian
sehingga $\int g_n\le\sum_{i\in\Bbb N}\int|f_i|$ untuk setiap $n$, sehingga
$g=\sup_{n\in\Bbb N}g_n$ terintegralkan, menurut teorema B. Levi sebagaimana dinyatakan dalam
123A.   Khususnya, $g$ berhingga hampir di mana-mana.   Sekarang, jika
$x\in X$ sedemikian sehingga $g(x)$ terdefinisi dan berhingga,
$\sum_{i\in J}|f_i(x)|\le g(x)$ untuk setiap $J\subseteq\Bbb N$ yang berhingga, sehingga
$\sum_{i\in\Bbb N}|f_i(x)|$ dan $\sum_{i\in\Bbb N}f_i(x)$ terdefinisi.
Dalam hal ini, tentu saja,
$\sum_{i\in\Bbb N}f_i(x)=\lim_{n\to\infty}\sum_{i=0}^nf_i(x)$.   Tetapi
$|\sum_{i=0}^nf_i|\leae g$ untuk setiap $n$, sehingga Teorema Konvergensi
Terdominasi Lebesgue menunjukkan bahwa

\Centerline{$\int\sum_{i\in\Bbb N}f_i
=\lim_{n\to\infty}\int\sum_{i=0}^nf_i
=\lim_{n\to\infty}\sum_{i=0}^n\int f_i
=\sum_{i\in\Bbb N}\int f_i$.}
}%end of proof of 226E

\exercises{
\leader{226X}{Latihan dasar $\pmb{>}$(a)}
%\sqheader 226Xa
Andaikan $I$ dan $J$ himpunan-himpunan dan
$\familyiI{a_i}$ suatu keluarga bilangan real yang dapat dijumlahkan.   (i) Tunjukkan bahwa
jika $f:J\to I$ injektif, maka $\family{j}{J}{a_{f(j)}}$ dapat dijumlahkan.
(ii) Tunjukkan bahwa jika $g:I\to J$ sembarang fungsi, maka
$\sum_{j\in J}\sum_{i\in g^{-1}[\{j\}]}a_i$ terdefinisi dan sama dengan
$\sum_{i\in I}a_i$.
%226A

\sqheader 226Xb Suatu {\bf fungsi tangga} pada
interval $[a,b]$ adalah fungsi $F$
sedemikian sehingga, untuk suatu pilihan $t_0,\ldots,t_n$ dengan $a=t_0\le\ldots\le t_n=b$,
$F$ konstan pada setiap interval $\ooint{t_{i-1},t_i}$.   Tunjukkan bahwa
$F:[a,b]\to\Bbb R$ merupakan fungsi saltus jika dan hanya jika untuk setiap $\epsilon>0$ terdapat
fungsi tangga $G:[a,b]\to\Bbb R$ sedemikian sehingga
$\Var_{[a,b]}(F-G)\le\epsilon$.
%226B

\spheader 226Xc Misalkan $F$, $G$ fungsi-fungsi bernilai real yang bervariasi
terbatas dan didefinisikan pada interval $[a,b]\subseteq\Bbb R$.   Tunjukkan bahwa, dalam
bahasa 226C,

\Centerline{$(F+G)_p=F_p+G_p$,\quad $(F+G)_c=F_c+G_c$,}

\Centerline{$(F+G)_{cs}=F_{cs}+G_{cs}$,
\quad$(F+G)_{ac}=F_{ac}+G_{ac}$.}
%226C

\sqheader 226Xd Misalkan $F$ suatu fungsi bernilai real yang bervariasi
terbatas pada interval $[a,b]\subseteq\Bbb R$.   Tunjukkan bahwa, dalam
bahasa 226C,

\Centerline{$\Var_{[a,b]}(F)
=\Var_{[a,b]}(F_p)+\Var_{[a,b]}(F_c)
=\Var_{[a,b]}(F_p)+\Var_{[a,b]}(F_{cs})+\Var_{[a,b]}(F_{ac})$.}
%226C

\spheader 226Xe Misalkan $F$ suatu fungsi bernilai real yang bervariasi
terbatas pada interval $[a,b]\subseteq\Bbb R$.   Tunjukkan bahwa $F$
kontinu mutlak jika dan hanya jika $\Var_{[a,b]}(F)=\int_a^b|F'|$.
%226Xd

\spheader 226Xf Tinjau fungsi $g$ dari 134I/226Cc.   Tunjukkan
bahwa $g^{-1}:[0,1]\to[0,1]$ terdiferensialkan hampir di mana-mana dalam
$[0,1]$, dan tentukan $\mu\{x:(g^{-1})'(x)\le a\}$ untuk setiap $a\in\Bbb R$.
%226C

\leader{226Y}{Latihan lanjutan (a)}
%\spheader 226Ya
Tunjukkan bahwa suatu himpunan $I$ terhitung jika dan hanya jika terdapat suatu keluarga
$\familyiI{a_i}$ dari bilangan-bilangan real tak nol yang dapat dijumlahkan.
%226A

\spheader 226Yb
Jelaskan modifikasi yang mungkin sesuai dalam deskripsi
dekomposisi Lebesgue dari suatu fungsi bervariasi terbatas
jika kita ingin meninjau fungsi pada interval terbuka atau
setengah terbuka, termasuk interval tak terbatas.
%226C

\spheader 226Yc Andaikan $F:[a,b]\to\Bbb R$ suatu fungsi bervariasi
terbatas, dan tetapkan $h(y)=\#(F^{-1}[\{y\}])$ untuk $y\in\Bbb R$.   Tunjukkan
bahwa $\int h=\Var_{[a,b]}(F_c)$, dengan $F_c$ `bagian kontinu'
dari $F$ sebagaimana didefinisikan dalam 226Ca/226Cd.
%226C

\spheader 226Yd Andaikan $a<b$ dalam $\Bbb R$, dan
$F:[a,b]\to\Bbb R$ suatu fungsi bervariasi terbatas;  misalkan $F_p$ bagian saltusnya.   Tunjukkan bahwa
$|F(b)-F(a)|\le\mu F[\,[a,b]\,]+\Var_{[a,b]}F_p$, dengan $\mu$ ukuran Lebesgue pada $\Bbb R$.
%226C
}%end of exercises


\endnotes{
\Notesheader{226} Dalam 232I dan 232Yb di bawah ini saya akan meninjau kembali gagasan-gagasan ini,
dengan mengaitkannya dengan suatu dekomposisi ukuran Lebesgue--Stieltjes
yang bersesuaian dengan suatu fungsi real tak-menurun, dan selanjutnya dengan ukuran-ukuran yang lebih
umum.   Seluruh pembahasan ini bersifat sampingan terhadap pokok perhatian
volume ini, tetapi menurut saya pembahasan tersebut mencerahkan, dan tentu saja merupakan bagian dari
pengetahuan dasar yang dianggap dimiliki oleh siapa pun yang bekerja dalam analisis real.
}%end of notes

\frnewpage
